\section{Introduction}
Non-Hermitian Hamiltonians with real spectra arise throughout mathematical physics and admit
a consistent Hilbert-space interpretation when they are self-adjoint with respect to a suitable
positive metric.  In finite dimensions, a diagonalizable matrix $H$ with real spectrum is
quasi-Hermitian: there exists $\eta=\eta^\dagger>0$ such that
$H^\dagger\eta=\eta H$, and $\rho=\eta^{1/2}$ transforms $H$ to the Hermitian representative
$h=\rho H\rho^{-1}$~\cite{Mostafazadeh2002I,Mostafazadeh2002II,Mostafazadeh2002III}.
The existence of such a metric and the associated similarity transformation are established
facts; the issue studied here is instead how the \emph{conditioning} of the metric interacts
with spectral inversion in the original reference norm.

The metric is generally nonunique.  Metric selection has therefore been studied using
additional principles.  Krej\v{c}i\v{r}\'ik, Lotoreichik and Znojil, for example, introduced a
minimally anisotropic metric through Hilbert--Schmidt-distance minimization and established
existence and uniqueness properties for that criterion~\cite{Krejcirik2018}.  Our objective is
different: we minimize the spectral condition number $\kappa_2(\eta)$.  This quantity controls
how severely the similarity map distorts Euclidean norms and therefore enters elementary
bounds that transport Hermitian inverse estimates back to the original representation.  The
condition-number problem should not be confused with the prior Hilbert--Schmidt optimization,
nor do we claim metric optimization itself as new.

The distinction becomes important near exceptional points.  Eigenvector coalescence is a
geometric feature of nonnormal spectral problems, and resolvent or perturbation behavior cannot
in general be inferred from eigenvalues alone~\cite{Trefethen2005}.  In pseudo-Hermitian
systems, the positive metric also becomes singular as a diagonalizable real-spectrum family
approaches a defective boundary.  It is tempting to read that singularity as evidence that
spectral inversion itself must become unstable.  The central question of this paper is whether
that implication is valid in the ordinary Euclidean representation.

Our answer is model-dependent but mathematically sharp.  We first derive general inequalities
relating the metric condition number to Euclidean inverse norms and perturbation radii.  A
two-level family shows that these inequalities can be exact.  We then introduce a coupled
three-level family $A_\varepsilon$ whose eigenvalues are $1$, $1+\varepsilon$ and $2$.
As $\varepsilon\downarrow0$, the first two eigenvectors coalesce and the limiting matrix becomes
defective at the \emph{nonzero} eigenvalue $1$.  Thus no eigenvalue approaches zero.  The family
provides a clean test of whether metric singularity is intrinsically tied to inversion
singularity.

The principal result is the sharp asymptotic characterization
\begin{equation}\label{eq:intro_main}
 \kappa_*(A_\varepsilon)
 :=\inf_{\eta>0,\;A_\varepsilon^\dagger\eta=\eta A_\varepsilon}
 \kappa_2(\eta)
 \sim \frac{8}{\varepsilon^2}.
\end{equation}
The lower bound uses the joint orthogonality constraints imposed by all three eigendirections,
while an explicit weighted-eigenvector construction attains the same leading constant.  The
third state therefore affects the optimal metric conditioning even though it remains spectrally
separated from the coalescing pair.  At the same time, $\|A_\varepsilon^{-1}\|_2$ and the exact
Euclidean distance to singularity remain uniformly bounded.  Hence the best metric-based
certificate collapses even when the inverse problem does not.  A separate perturbation theorem
shows that normalized solutions remain uniformly stable.

We finally embed the same mechanism into a coupled $N$-level bidiagonal family.  The lower
bound $\kappa_*\ge 8/\varepsilon^2$ survives for every $N$, whereas the inverse norm remains
bounded uniformly in both $N$ and $\varepsilon$.  A standard Hermitian dilation is recorded as
a computational implication, not as a new quantum algorithm: its condition number depends on
the singular values of the original matrix and stays bounded for our growing family.  This
makes explicit why a divergent pseudo-Hermitian metric condition number is not, by itself, a
complexity statement about quantum linear-system inversion~\cite{Gilyen2019,Berry2026}.

The paper is organized as follows.  Section~2 fixes the metric geometry and norm conventions.
Section~3 derives the general inversion bounds.  Section~4 gives the sharp two-level benchmark.
Section~5 develops the coupled three-level counterexample and proves~\eqref{eq:intro_main}.
Section~6 extends the separation to growing dimension and states the metric-free Hermitian
baseline.  Section~7 reports reproducible numerical checks, and Sections~8--9 discuss the
scope and conclusions.
